\subsection{Wall jump a odraz od stěny}

Wall jump se provádí pouze když hráč stiskne skok, nachází se na stěně a není na zemi. Skript detekuje směr stěny pomocí \texttt{DetectWallDirection()} a aplikuje sílu opačným směrem. Vertikální složka skoku je pevně nastavena na \texttt{wallJumpPowerY}, zatímco horizontální složka závisí na typu skoku.

Implementace využívá \texttt{Rigidbody2D} pro aplikaci síly metodou \texttt{AddForce()}. Vektory se nastavují podle směru stěny: pokud je stěna vpravo, horizontální složka je záporná, pokud je stěna vlevo, horizontální složka je kladná.

Tento mechanismus umožňuje hráči systematicky se odrážet mezi dvěma stěnami a dosahovat vyšších částí úrovně, kam by se normálním skokem nedostal. Wall jump je základní technikou v mnoha platformovkách a přidává hloubku do navigace herním světem.

\begin{lstlisting}[language=csharp,caption={Listing 11: Metoda HandleWallJump}]
void HandleWallJump()
{
    if (!pc.jumpPressedThisFrame || lastOnWallTime <= 0 || pc.grounded) return;

    isWallJumping = true;
    int jumpDir = -wallDir;

    bool isWallToWall = Physics2D.Raycast(wallCheck.position, 
        Vector2.right * jumpDir, 1f, wallLayer).collider == null;

    float xForce = isWallToWall ? data.wallToWallJumpForceX : data.singleWallJumpForceX;
    rb.linearVelocity = new Vector2(jumpDir * xForce, data.wallJumpPowerY);

    if (pc.isFacingRight != (jumpDir > 0))
        pc.CheckDirectionToFace(jumpDir > 0);

    lastOnWallTime = 0;
}
\end{lstlisting}
